% Čest: logika-prvogo-reda
% Glava: teorija-modelov
% Odděl: izomorfizm

\documentclass[../../../include/open-logic-section]{subfiles}

\begin{document}

\olfileid[isv]{mod}{bas}{iso}
\olsection{Izomorfne struktury}

!!{structure}s logiky prvogo reda mogut byti podobne na dva načina.
Jedin način podobnosti jest, že v njih su istinne te same
!!{sentence}s. Takove !!{structure}s nazyvajemo
\emph{elementarno ekvivalentnymi}. No struktury mogut se mnogo
razlikovati i vseže iměti istinne te same !!{sentence}s---
na priměr, jedna može byti !!{enumerable}, a druga ne. To jest zato, že
sučstvuje mnogo svojstv !!a{structure}, ktore ne mogut byti izražene
v jezykah prvogo reda, ili zato, že jezyk ne jest dostatočno bogaty,
ili iz-za fundamentalnyh ograničenij logiky prvogo reda, kako teorema
L\"owenheima--Skolema. Drugy, strožši način podobnosti
!!{structure}s jest, že one sųt v osnově jednake: razlikujut se jedino
objektami, iz ktoryh sųt sostavljene, a ne svojimi strukturnymi
svojstvami. To možno precizno izraziti pojmom \emph{izomorfizma}.

\begin{defn}
\ollabel{defn:elem-equiv} 
Dane sųt dve !!{structure}s $\Struct{M}$ i $\Struct M'$ jednogo
!!{language}~$\Lang{L}$. Govorimo, že $\Struct{M}$ jest
\emph{elementarno ekvivalentna} strukturi $\Struct M'$, i pišemo
$\Struct{M} \equiv \Struct M'$, togda i samo togda, kogda dla vsakogo
!!{sentence}~$!A$ jezyka~$\Lang{L}$,
$\Sat{M}{!A}$ togda i samo togda, kogda $\Sat{M'}{!A}$.
\end{defn}

\begin{defn}
\ollabel{defn:isomorphism} Dane sųt dve !!{structure}s $\Struct{M}$ i
$\Struct M'$ jednogo !!{language}~$\Lang L$. Govorimo, že
$\Struct{M}$ jest \emph{izomorfna} strukturi $\Struct M'$, i pišemo $\Struct{M}
\simeq \Struct M'$, togda i samo togda, kogda sučstvuje funkcija $h
\colon \Domain{M} \to \Domain{M'}$ taka, že:
\begin{enumerate}
\item $h$ jest !!{injective}: ako $h(x) =
  h(y)$, togda $x = y$; 
\item $h$ jest !!{surjective}: dla vsakogo $y \in \Domain{M'}$
  sučstvuje $x \in \Domain{M}$ tako, že $h(x) = y$;
\item \ollabel{defn:iso-const}dla vsakoj !!{constant} $c$:
  $h(\Assign{c}{M}) = \Assign{c}{M'}$;
\item \ollabel{defn:iso-pred}dla vsakogo $n$-městnogo !!{predicate}~$P$:
  \[
  \tuple{a_1, \dots, a_n}\in \Assign{P}{M} \quad\text{togda i samo togda, kogda}\quad
  \tuple{h(a_1), \dots, h(a_n)} \in \Assign{P}{M'};
  \]
\item \ollabel{defn:iso-func}dla vsakoj $n$-městnoj !!{function} $f$:
  \[
  h(\Assign{f}{M}(a_1, \dots, a_n)) =
  \Assign{f}{M'}(h(a_1), \dots, h(a_n)).
  \]
\end{enumerate}
\end{defn}

\begin{thm}
\ollabel{thm:isom}
Ako $\Struct{M} \iso \Struct M'$, togda $\Struct{M} \elemequiv
\Struct{M'}$.
\end{thm}

\begin{proof}
Nehaj $h$ bude izomorfizmom struktury $\Struct{M}$ na strukturu
$\Struct M'$. Dla vsakogo prisvajanja~$s$, $h \circ s$ jest
kompozicija funkcij $h$ i $s$, t.j. prisvajanje v $\Struct{M'}$ tako,
že $(h \circ s)(x) = h(s(x))$. Indukcijeju po složenosti $t$ i $!A$
možno dokazati slědujuće silnějše tvrđenja:
\begin{enumerate}
  \item[a.] $h(\Value{t}{M}[s]) = \Value{t}{M'}[h\circ s]$.
  \item[b.] $\Sat{M}{!A}[s]$ togda i samo togda, kogda $\Sat{M'}{!A}[h \circ s]$.
\end{enumerate}
Prvo tvrđenje dokazyvaje se indukcijeju po složenosti~$t$.
\begin{enumerate}
\item Ako $t \ident c$, togda $\Value{c}{M}[s] = \Assign{c}{M}$ i
  $\Value{c}{M'}[h \circ s] = \Assign{c}{M'}$. Zato
  $h(\Value{t}{M}[s]) = h(\Assign{c}{M}) = \Assign{c}{M'}$ (po
  \olref{defn:iso-const} iz \olref{defn:isomorphism}) $=
  \Value{t}{M'}[h \circ s]$.
\item Ako $t \ident x$, togda $\Value{x}{M}[s] = s(x)$ i
  $\Value{x}{M'}[h \circ s] = h(s(x))$. Zato $h(\Value{x}{M}[s]) =
  h(s(x)) = \Value{x}{M'}[h \circ s]$.
\item Ako $t \ident f(t_1, \dots, t_n)$, togda
  \begin{align*}
    \Value{t}{M}[s] & = \Assign{f}{M}(\Value{t_1}{M}[s], \dots, \Value{t_n}{M}[s]) \quad\text{i}\\
  \Value{t}{M'}[h \circ s] & = \Assign{f}{M'}(\Value{t_1}{M'}[h \circ
    s], \dots, \Value{t_n}{M'}[h \circ s]).
  \end{align*}
  Indukcijska hipoteza govori, že dla vsakogo $i$, $h(\Value{t_i}{M}[s])
  = \Value{t_i}{M'}[h\circ s]$. Zato
  \begin{align}
    h(\Value{t}{M}[s]) 
    & = h(\Assign{f}{M}(\Value{t_1}{M}[s], \dots, \Value{t_n}{M}[s]))\notag\\
    & = \Assign{f}{M'}(h(\Value{t_1}{M}[s]), \dots,
    h(\Value{t_n}{M}[s])) \ollabel{iso-1}\\
    & = \Assign{f}{M'}(\Value{t_1}{M'}[h \circ s], \dots,
    \Value{t_n}{M'}[h \circ s]) \ollabel{iso-2}\\
    & = \Value{t}{M'}[h\circ s] \notag
  \end{align}
  Tut \olref{iso-1} slěduje iz \olref{defn:iso-func} definicije
  \olref{defn:isomorphism}, a \olref{iso-2} iz indukcijske hipotezy.
\end{enumerate}
Dokaz česti (b) ostavljajemo kako zadanje.

Ako $!A$ jest izrečenje, prisvajanja~$s$ i $h \circ s$ ne imajut
vliva, i imajemo $\Sat{M}{!A}$ togda i samo togda, kogda
$\Sat{M'}{!A}$.
\end{proof}

\begin{prob}
Izvedite podrobny dokaz česti (b) iz
\olref[mod][bas][iso]{thm:isom}. Ne zabudte ukazati, kde jest
upotrěbljeno vsako iz pěti svojstv, ktore harakterizujut izomorfizmy
iz \olref[mod][bas][iso]{defn:isomorphism}.
\end{prob}

\begin{defn}
\emph{Avtomorfizm} struktury $\Struct{M}$ jest izomorfizm struktury
$\Struct{M}$ na samu sebe.
\end{defn}

\begin{prob}
Dokažite, že dla vsakoj struktury $\Struct{M}$, ako $X$ jest
definovateljno podmnožstvo struktury $\Struct{M}$, a $h$ jest
avtomorfizm struktury $\Struct{M}$, togda $X = \Setabs{h(x)}{x \in X}$
(t.j. $X$ ostavaje neizměnno pri $h$).
\end{prob}


\end{document}
